\SourcePage{791}{794}
\renewcommand{\thesection}{\alph{section}}
\setcounter{section}{5}
\section{Evaluation of integrals containing confluent hypergeometric functions}

Consider an integral of the form
\begin{equation}
 J_{\alpha\gamma}^{\nu}=\int_0^\infty
 e^{-\lambda z}z^\nu F(\alpha,\gamma,kz)\,dz.
 \tag{f.1}
\end{equation}
It is assumed to converge. This requires
$\operatorname{Re}\nu>-1$ and $\operatorname{Re}\lambda>|\operatorname{Re}k|$;
if $\alpha$ is a negative integer,

\SourcePage{792}{795}
the second condition may instead be replaced by
$\operatorname{Re}\lambda>0$. Use the integral representation (d.9) for
$F(\alpha,\gamma,kz)$ and carry out the integration with respect to $z$ under
the contour-integral sign. This gives
\[
 \begin{split}
 J_{\alpha\gamma}^{\nu}
 &=-\frac1{2\pi i}\frac{\Gamma(1-\alpha)\Gamma(\gamma)}
 {\Gamma(\gamma-\alpha)}\lambda^{-\nu-1}\Gamma(\nu+1)\\
 &\hspace{1em}\times\oint_{C'}(-t)^{\alpha-1}(1-t)^{\gamma-\alpha-1}
 \left(1-\frac{k}{\lambda}t\right)^{-\nu-1}dt.
 \end{split}
\]
Taking (e.3) into account, we finally find
\begin{equation}
 J_{\alpha\gamma}^{\nu}=\Gamma(\nu+1)\lambda^{-\nu-1}
 F\left(\alpha,\nu+1,\gamma,\frac{k}{\lambda}\right).
 \tag{f.2}
\end{equation}
In the cases where $F(\alpha,\nu+1,\gamma,k/\lambda)$ reduces to
polynomials, the corresponding expressions for
$J_{\alpha\gamma}^{\nu}$ are in elementary functions:
\begin{equation}
 J_{\alpha\gamma}^{\gamma+n-1}=(-1)^n\Gamma(\gamma)
 \frac{d^n}{d\lambda^n}
 \left[\lambda^{\alpha-1}(\lambda-k)^{-\alpha}\right],
 \tag{f.3}
\end{equation}
\begin{equation}
 J_{-n,\gamma}^{\nu}=(-1)^n
 \frac{\Gamma(\nu+1)(\lambda-k)^{\gamma+n-\nu-1}}
 {\gamma(\gamma+1)\cdots(\gamma+n-1)}
 \frac{d^n}{d\lambda^n}
 \left[\lambda^{-\nu-1}(\lambda-k)^{\nu-\gamma+1}\right],
 \tag{f.4}
\end{equation}
\begin{equation}
 \begin{split}
 J_{\alpha m}^{n}
 &=\frac{(-1)^{m-n}}
 {k^{m-1}(1-\alpha)(2-\alpha)\cdots(m-1-\alpha)}\\
&\hspace{1em}\times\Biggl\{-(m-1)!\frac{d^n}{d\lambda^n}
 \left[\lambda^{\alpha-1}(\lambda-k)^{m-\alpha-1}\right]\\
&\hspace{1em}+n!(m-n-1)\cdots(m-1)
 \lambda^{\alpha-n-1}(\lambda-k)^{-1+m-n-\alpha}\\
&\hspace{2em}\times\frac{d^{m-n-2}}{d\lambda^{m-n-2}}
 \left[\lambda^{m-\alpha-1}(\lambda-k)^{\alpha-1}\right]\Biggr\},
 \end{split}
 \tag{f.5}
\end{equation}
where $m,n$ are integers and $0\le n\le m-2$.

Next, evaluate the integral
\begin{equation}
 J_\nu=\int_0^\infty e^{-kz}z^{\nu-1}
 [F(-n,\gamma,kz)]^2\,dz
 \tag{f.6}
\end{equation}
($n$ is a positive integer, $\operatorname{Re}\nu>0$). We begin with the
more general integral having $e^{-\lambda z}$ rather than $e^{-kz}$ in its
integrand. Write one of the functions $F(-n,\gamma,kz)$ as the integral (d.9),
after which integration with respect to $z$ by

\SourcePage{793}{796}
formula (f.3) gives
\[
 \begin{split}
 &\int_0^\infty e^{-\lambda z}z^{\nu-1}[F(-n,\gamma,kz)]^2\,dz\\
 &=-\frac1{2\pi i}(-1)^n
 \frac{\Gamma(1+n)\Gamma^2(\gamma)\Gamma(\nu)}
 {\Gamma^2(\gamma+n)}\\
 &\hspace{1em}\times\oint_{C'}(\lambda-kt-k)^{\gamma+n-\nu}
 (-t)^{-n-1}(1-t)^{\gamma+n-1}\\
 &\hspace{3em}\times\frac{d^n}{d\lambda^n}
 \left[(\lambda-kt)^{-\nu}(\lambda-kt-k)^{\nu-\gamma}\right]dt.
 \end{split}
\]
The derivative of order $n$ with respect to $\lambda$ can plainly be
replaced by an expression in terms of a derivative of the same order with
respect to $t$. Having done this, put $\lambda=k$, thereby returning to the
integral $J_\nu$:
\[
 \begin{split}
 J_\nu&=-\frac1{2\pi i}
 \frac{\Gamma(n+1)\Gamma(\nu)\Gamma^2(\gamma)}
 {\Gamma^2(\gamma+n)k^\nu}\\
 &\hspace{1em}\times\oint_{C'}(-t)^{\gamma-\nu-1}(1-t)^{\gamma+n-1}
 \frac{d^n}{dt^n}\left[(1-t)^{-\nu}(-t)^{\nu-\gamma}\right]dt.
 \end{split}
\]
Integrating by parts $n$ times transfers the operation $(d/dt)^n$ to
$(-t)^{\gamma-\nu-1}(1-t)^{\gamma+n-1}$; expand the derivative by Leibniz's
formula. The result is a sum of integrals, each reducible to the familiar
Euler integral. The final expression for the integral sought is
\begin{equation}
 \begin{split}
 J_\nu&=\frac{\Gamma(\nu)n!}
 {k^\nu\gamma(\gamma+1)\cdots(\gamma+n-1)}\\
&\hspace{1em}\times\left\{1+\sum_{s=0}^{n-1}
 \frac{n(n-1)\cdots(n-s)(\gamma-\nu-s-1)
 (\gamma-\nu-s)\cdots(\gamma-\nu+s)}
 {[(s+1)!]^2\gamma(\gamma+1)\cdots(\gamma+s)}\right\}.
 \end{split}
 \tag{f.7}
\end{equation}
It is readily seen that the integrals $J_\nu$ satisfy the following relation
($p$ is an integer):
\begin{equation}
 J_{\gamma+p}=\frac{(\gamma-p-1)(\gamma-p)\cdots(\gamma+p-1)}
 {k^{2p+1}}J_{\gamma-1-p}.
 \tag{f.8}
\end{equation}
In the same way one evaluates the integral
\begin{equation}
 J=\int_0^\infty e^{-\lambda z}z^{\gamma-1}
 F(\alpha,\gamma,kz)F(\alpha',\gamma,k'z)\,dz.
 \tag{f.9}
\end{equation}

\SourcePage{794}{797}
Represent $F(\alpha',\gamma,k'z)$ by the integral (d.9). After integrating
with respect to $z$ by formula (f.3), with $n=0$, we obtain
\[
 \begin{split}
 J&=-\frac1{2\pi i}\frac{\Gamma(1-\alpha')\Gamma^2(\gamma)}
 {\Gamma(\gamma-\alpha')}\oint_{C'}(-t)^{\alpha'-1}
 (1-t)^{\gamma-\alpha'-1}\\
 &\hspace{3em}\times(\lambda-k't)^{\alpha-\gamma}
 (\lambda-k't-k)^{-\alpha}\,dt.
 \end{split}
\]
The substitution $t\to\lambda t/(k't+\lambda-k')$ reduces this integral to
the form (e.3), giving
\begin{equation}
 J=\Gamma(\gamma)\lambda^{\alpha+\alpha'-\gamma}
 (\lambda-k)^{-\alpha}(\lambda-k')^{-\alpha'}
 F\left(\alpha,\alpha',\gamma,
 \frac{kk'}{(\lambda-k)(\lambda-k')}\right).
 \tag{f.10}
\end{equation}
If $\alpha$ (or $\alpha'$) is a negative integer, $\alpha=-n$, relation
(e.7) can be used to rewrite this expression as
\begin{equation}
 \begin{split}
 J&=\frac{\Gamma^2(\gamma)\Gamma(\gamma+n-\alpha')}
 {\Gamma(\gamma+n)\Gamma(\gamma-\alpha')}
 \lambda^{-n+\alpha'-\gamma}(\lambda-k)^n
 (\lambda-k')^{-\alpha'}\\
 &\hspace{2em}\times F\left(-n,\alpha',-n+\alpha'+1-\gamma,
 \frac{\lambda(\lambda-k-k')}{(\lambda-k)(\lambda-k')}\right).
 \end{split}
 \tag{f.11}
\end{equation}

Finally, consider integrals of the form
\begin{equation}
 J_\gamma^{sp}(\alpha,\alpha')=\int_0^\infty
 \exp\left(-\frac{k+k'}2z\right)z^{\gamma-1+s}
 F(\alpha,\gamma,kz)F(\alpha',\gamma-p,k'z)\,dz.
 \tag{f.12}
\end{equation}
The parameters are assumed to have values for which the integral converges
absolutely; $s,p$ are positive integers. By (f.10), the simplest of these
integrals, $J_\gamma^{00}(\alpha,\alpha')$, is
\begin{equation}
 \begin{split}
 J_\gamma^{00}(\alpha,\alpha')
 &=2^\gamma\Gamma(\gamma)(k+k')^{\alpha+\alpha'-\gamma}
 (k'-k)^{-\alpha}(k-k')^{-\alpha'}\\
 &\hspace{2em}\times F\left(\alpha,\alpha',\gamma,
 -\frac{4kk'}{(k'-k)^2}\right),
 \end{split}
 \tag{f.13}
\end{equation}
and if $\alpha$ (or $\alpha'$) is a negative integer ($\alpha=-n$), then,

\SourcePage{795}{798}
by (f.11), it may also be written as
\begin{equation}
 \begin{split}
 J_\gamma^{00}(-n,\alpha')
 &=2^\gamma\frac{\Gamma(\gamma)(\gamma-\alpha')
 (\gamma-\alpha'+1)\cdots(\gamma-\alpha'+n-1)}
 {\gamma(\gamma+1)\cdots(\gamma+n-1)}\\
 &\hspace{1em}\times(-1)^n(k+k')^{-n+\alpha'-\gamma}
 (k-k')^{n-\alpha'}\\
 &\hspace{1em}\times F\left[-n,\alpha',\alpha'+1-n-\gamma,
 \left(\frac{k+k'}{k-k'}\right)^2\right].
 \end{split}
 \tag{f.14}
\end{equation}
A general formula for $J_\gamma^{sp}(\alpha,\alpha')$ can be derived, but it
is too complicated for convenient use. It is preferable to use recurrence
relations that reduce $J_\gamma^{sp}(\alpha,\alpha')$ to the integral with
$s=p=0$\footnote{See W.~Gordon // Ann. d. Phys. 1929. V.~2. P.~1031.}. The
formula
\begin{equation}
 J_\gamma^{sp}(\alpha,\alpha')=\frac{\gamma-1}{k}
 \left\{J_{\gamma-1}^{s,p-1}(\alpha,\alpha')
 -J_{\gamma-1}^{s,p-1}(\alpha-1,\alpha')\right\}
 \tag{f.15}
\end{equation}
makes it possible to reduce $J_\gamma^{sp}(\alpha,\alpha')$ to an integral
with $p=0$. Thereafter the formula
\begin{equation}
 \begin{split}
 J_\gamma^{s+1,0}(\alpha,\alpha')
 &=\frac4{k^2-k'^2}\Biggl\{
 \left[\frac\gamma2(k-k')-k\alpha+k'\alpha'-k's\right]
 J_\gamma^{s0}(\alpha,\alpha')\\
 &\hspace{2em}+s(\gamma-1+s-2\alpha')
 J_\gamma^{s-1,0}(\alpha,\alpha')
 +2\alpha'sJ_\gamma^{s-1,0}(\alpha,\alpha'+1)\Biggr\}
 \end{split}
 \tag{f.16}
\end{equation}
completes the reduction to the integral with $s=p=0$.
